% sink_results.tex -- results prose. All numbers come from \newcommand{\NSettings}{11}
\newcommand{\NModels}{5}
\newcommand{\NCellsTotal}{568,028}
\newcommand{\NSinkSettings}{10}
\newcommand{\PctConedRange}{36--88\%}
\newcommand{\SinkMassRange}{0.42 to 0.69}
\newcommand{\NoSinkMass}{0.02}
\newcommand{\RhoMin}{0.999}
\newcommand{\RhoNoSink}{0.967}
\newcommand{\RhoDeltaHone}{0.33 to 0.89}
\newcommand{\NViolations}{0}
\newcommand{\MaxSeedSD}{0.008}
\newcommand{\NComp}{10}
\newcommand{\NTOneEquiv}{4}
\newcommand{\TOneRange}{-0.015 to +0.031}
\newcommand{\TTwoRange}{+0.004 to +0.030}
\newcommand{\NTTwoEquiv}{3}
\newcommand{\TThreeTQA}{+0.032 to +0.045}
\newcommand{\TThreeOnp}{+0.002 to +0.029}
\newcommand{\TThreeHE}{+0.052 to +0.055}
\newcommand{\NTFourEquiv}{8}
\newcommand{\TFourRange}{-0.008 to +0.026}
\newcommand{\PHTopoRange}{-0.001 to +0.011}
\newcommand{\NPHTopoEquiv}{9}
\newcommand{\NPHTopoTotal}{9}
\newcommand{\PHEntRange}{-0.016 to +0.036}
\newcommand{\NIZeroEquiv}{10}
\newcommand{\NIZeroTotal}{10}
\newcommand{\IZeroRange}{0.000 to +0.001}
\newcommand{\NIDeflEquiv}{10}
\newcommand{\IDeflRange}{-0.001 to +0.002}
\newcommand{\LexTQA}{0.82}
\newcommand{\LexOnp}{0.60 to 0.74}
\newcommand{\LexHE}{0.98}
\newcommand{\LogpTQA}{0.58 to 0.64}
\newcommand{\LogpOnp}{0.72 to 0.88}
\newcommand{\ZeroDTQA}{0.58 to 0.75}
\newcommand{\ZeroDOnp}{0.71 to 0.73}
\newcommand{\ZeroDHE}{0.79 to 0.82}
\newcommand{\NonTopoAll}{0.84 to 0.94}
\newcommand{\HiddenAll}{0.83 to 0.93}
\newcommand{\LookbackAll}{0.82 to 0.90}
\newcommand{\HiddenHE}{1.00}
\newcommand{\SplitConed}{-0.001 to +0.004}
\newcommand{\SplitOpen}{+0.003 to +0.029}
\newcommand{\OnpAccRange}{26\% to 56\%}
,
% which sinktda/report.py regenerates from sinktda_results/.

\newcommand{\RESULTSREDUCTION}{%
Table~\ref{tab:theory} checks Theorem~\ref{thm:cone} on every layer-level graph of every example: \NCellsTotal{} graphs in total. There are \NViolations{} violations of any bound. In the \NSinkSettings{} settings whose models place \SinkMassRange{} of attention on token~0, \PctConedRange{} of graphs are \emph{exactly} coned ($\defect_0=0$), and on every one of them $H_1$ is empty, as prediction C2 requires. Across graphs, $\defect_0$ and $1$D total persistence have Spearman correlation \RhoDeltaHone{} (C3), the floor being the \RhoDeltaHoneNearN{} settings above \RhoDeltaHoneNearCut{} coned, where both quantities are near-constant (\RhoDeltaHoneNearHOneZero{} of graphs have no $1$D bars) and the correlation is degenerate rather than weak; elsewhere it is at least \RhoDeltaHoneRest{}. Both $P_0$ and the star weight $\stw_0=(N{-}1)(1-\bar m)$ scale with sequence length, so their raw per-layer rank correlation (median $\ge$\RhoRawMin{}) partly reflects length. After dividing both by $N{-}1$, the median per-layer rank correlation between $P_0/(N{-}1)$ and $1-\bar m$ is still at least \RhoMin{} in every sink setting (C1). In most layers, then, length-normalized $0$D total persistence is, up to rank, one minus the mean attention paid to the first token. The weakest single layer per setting reaches only \RhoLayerMinRange{}.

Figure~\ref{fig:layers} shows where coning happens: in the middle of the network, where sink attention is high. The first few layers are never sink-dominated, and in some models (TinyLlama, the Qwen models) so is a block of late layers; these layers are rarely coned. There, the length-normalized correlation can drop, as Theorem~\ref{thm:cone}(b) allows once $(N{-}1)\defect_0$ is comparable to the spread of $\stw_0$ across examples.

\paragraph{Models without a first-token sink as a contrast.} Our study includes \NNoSinkModels{} models that place almost no attention on token~0 --- \NoSinkModels{}, at \NoSinkMass{} --- and none of their graphs is coned at $s=0$, so Theorem~\ref{thm:cone} at the first token gives no useful guarantee for them. They are not alike. Allowing any apex, \NoSinkConedAnyMin{} of \NoSinkConedAnyMinModel{}'s graphs are coned --- no vertex cones any of them --- against \NoSinkConedAnyMax{} of \NoSinkConedAnyMaxModel{}'s, whose attention is strongly concentrated, just not on token~0 (\S\ref{sec:toha}): one model has no cone structure to find, the other is one where $s=0$ is the wrong apex. Their $H_1$ is non-empty on \NoSinkHOneNonempty{} of graphs, and although the raw $\rho(P_0,\stw_0)$ is \RhoNoSinkRaw{}, this is almost entirely shared length: the length-normalized correlation is \RhoNoSink{}. Both differ from the sink-dominated models in weights, data and template, so this is a contrast, not a controlled intervention; the within-model intervention is deflation (C4, below). The reduction is a property of sink-dominated attention, not of the pipeline.
}

\newcommand{\RESULTSBEYOND}{%
If $0$D persistence reduced \emph{exactly} to the star, a classifier on the two sink features would match one on the two $0$D features in every layer. Table~\ref{tab:tests} (``Sink vs.\ 0D'') shows the two banks differ by \TOneRange{} AUC across the \NComp{} sink settings evaluated, and are statistically equivalent within $\pm0.015$ in \NTOneEquiv{} of them (the TOST power of this comparison is modest, so non-equivalence is not evidence of a difference). Adding $0$D to the sink bank gains \TTwoRange{}.

\paragraph{The residual lives exactly where the theory says it can.} We split layers without using labels: \emph{coned} layers ($\ge90\%$ of graphs coned) and \emph{open} layers ($<50\%$). Adding the coned layers' $0$D features to \textsc{Sink} changes AUC by \SplitConed{}, whereas the open layers account for \SplitOpen{} (Appendix~\ref{app:tables}). Almost everything $0$D adds beyond sink attention comes from the layers where the sink does not dominate.

\paragraph{Per head, the reduction is tighter.} Per-head sink mass and per-head $0$D are far stronger than their head-averaged versions (Table~\ref{tab:auc}), but adding per-head MST total ($P_{0,h}$) to per-head sink mass changes AUC by only \PHTopoRange{}, which is equivalent within $\pm0.015$ in \NPHTopoEquiv{} of the \NPHTopoTotal{} sink settings with per-head features.

\paragraph{The coning defect is a signal, but not a new one.} Theorem~\ref{thm:cone} singles out $\defect_0$ as the only place where topology can depart from the sink, so we also test it as a detector. Head-averaged, it is weak (AUC \DefLayerAuc{}). Per head it is strong (\DefHeadAuc{}), and stacked onto per-head sink mass it adds \DefHSinkRange{} AUC (95\% interval above $0$ in \DefHSinkSig{} of \DefTotal{} settings). Stacked onto the non-topological probes, however, it adds only \DefHNTRange{} (above $0$ in \DefHNTSig{}, equivalent within $\pm0.015$ in \DefHNTEquiv{} of \DefTotal{}).

\paragraph{Deleting the sink exposes some, redundant, structure.} Deflated PH (C4) adds \TThreeTQA{} over \textsc{Sink} on TruthfulQA (95\% interval excludes $0$ in \NDeflTQASig{} of \NDeflTQATotal{} sink settings) but only \TThreeOnp{} on the on-policy benchmark (\NDeflOnpSig{} of \NDeflOnpTotal{}). Answer-only PH behaves similarly (Appendix~\ref{app:tables}). So attention graphs carry some non-sink geometry, most clearly on teacher-forced text. As \S\ref{sec:detection} shows, this geometry is already available to non-topological probes.
}

\newcommand{\RESULTSONEDIM}{%
Under i.i.d.\ weights, $1$D persistence grows faster than linearly over the lengths that matter here and linearly asymptotically (Proposition~\ref{prop:morse}; Figure~\ref{fig:synthetic}, left). Real attention behaves very differently: most graphs have \emph{no} $1$D bars at all (Table~\ref{tab:theory}), and the bars that exist are confined to graphs with a non-zero coning defect. Figure~\ref{fig:synthetic} (middle) reproduces this transition in a controlled setting: raising the sink's logit bias moves random causal attention from the i.i.d.-like regime to exact coning, with $H_1$ emptying in lockstep.

This explains a recurring empirical observation, that $1$D features add little to $0$D. At full sample size, adding length-normalized $1$D to $0$D changes AUC by \TFourRange{} and is equivalent within $\pm0.015$ in \NTFourEquiv{} of \NComp{} sink settings; the exceptions are \TFourExceptions{}; of these, adding $1$D significantly improves AUC in \TFourHelps{}, while the remaining one has a negative point estimate whose interval is too wide to certify equivalence. The TruthfulQA Qwen2.5-3B gain is also the result that is most sensitive to inference precision (Appendix~\ref{app:impl}). The mechanism is not that loops are uninformative in principle. In a strongly sinked model, a loop can only exist where some token pays another token more attention than one of the two pays to the sink, and its lifetime is bounded by that excess. Figure~\ref{fig:synthetic} (right) makes this concrete. A planted $6$-cycle that $H_1$ detects easily under a moderate sink ($b{=}6$, $\gamma{=}8$: AUC \PlantModerate{}) is invisible under a strong one ($b{=}8$, $\gamma\le4$: AUC $\le$\PlantStrong{}) until it is strong enough to break the cone. Without any sink ($b{=}0$) $H_1$ also misses it (AUC $\le$\PlantNone{}), because random attention already has many background cycles; $1$D persistence is informative only in between.
}

\newcommand{\RESULTSDETECTION}{%
Table~\ref{tab:auc} places attention-graph topology among standard detectors. Three patterns hold across benchmarks.

\emph{Topology is dominated by cheap non-topological probes.} On TruthfulQA and the on-policy benchmark, a mean-pooled hidden-state probe (\HiddenAll{}) and per-head Lookback ratios (\LookbackAll{}) outperform every topological bank, and their combination with log-probability, LLM-Check, and row statistics reaches \NonTopoAll{}. Adding $0$D PH to that combination changes AUC by \IZeroRange{} (equivalent within $\pm0.015$ in \NIZeroEquiv{} of the \NIZeroTotal{} settings with per-head features), and adding sink-deflated PH changes it by \IDeflRange{} (equivalent in \NIDeflEquiv{} of \NIZeroTotal{}). Early fusion could hide a small bank inside a ${\sim}2{,}000$-dimensional one, so we repeat the test with late fusion. Stacking the $0$D score onto the non-topological score changes AUC by \LFZeroNTRange{} (equivalent in \LFZeroNTEquiv{} of \LFZeroNTTotal{} settings; the 95\% interval lies above $0$ in \LFZeroNTSig{} of them), and stacking the deflated-PH score changes it by \LFDeflNTRange{} (equivalent in \LFDeflNTEquiv{} of \LFDeflNTTotal{}).

\emph{Which cheap probe wins depends on whether text is teacher-forced.} On TruthfulQA, token log-probability is weak (\LogpTQA{}), consistent with the benchmark's design around imitative falsehoods, and $0$D PH (\ZeroDTQA{}) scores higher for \NLogpBeatsZeroDTQA{} of \NLogpZeroDTQATotal{} models, significantly so for \NLogpZeroDTQA{}. On-policy, log-probability is strong (\LogpOnp{}) and significantly beats $0$D (\ZeroDOnp{}) for \NLogpOnpSig{} of \NLogpOnpTotal{} models. Conclusions about topology versus confidence drawn from a single teacher-forced benchmark do not transfer.

\emph{Per-head resolution matters more than topology.} Per-head features of every kind are far stronger than head-averaged ones, and per-head entropy lands within \PHEntRange{} AUC of per-head $0$D (Appendix~\ref{app:tables}).
}

\newcommand{\RESULTSARTIFACTS}{%
A TF-IDF model on the answer text alone reaches \LexTQA{} AUC on TruthfulQA and \LexHE{} on HaluEval, but only \LexOnp{} on-policy: teacher-forced benchmarks reward recognizing how the dataset's answers were \emph{written} \citep{hussain2026parallax,janiak2025illusion}. We draw no conclusions from HaluEval, where lexical separability is near perfect.
}

\newcommand{\RESULTSDISCUSSION}{%
\paragraph{What attention-graph topology measures.} In decoder-only models, Vietoris--Rips persistence of an attention graph, and every stable summary of it, is controlled by one scalar per graph: the coning defect of the first token. The same is true of TOHA's prompt--response divergence, with the prompt playing the role of the apex. Where it is zero, as in most middle-layer graphs of every sink-dominated model we tested, $0$D total persistence equals $(N{-}1)(1-\bar m)$, a function of length and sink attention alone, and higher homology is empty. Where it is positive, Theorem~\ref{thm:cone} bounds how far topology can depart from the sink; in our data that is mainly the first layers. Practical recommendations follow in Appendix~\ref{app:recs}.
}

\newcommand{\MISTRALAPPENDIX}{%
Mistral-7B-Instruct-v0.2 is the most sink-dominated model in this study, and the reduction is correspondingly tightest there. Under its \texttt{[INST]} format it places \MistralSinkMass{} of its attention on token~0 --- more than any smaller model --- and \MistralConed{} of its \MistralCells{} layer graphs are \emph{exactly} coned ($\defect_0=0$), with \MistralViol{} violations of any bound in Theorem~\ref{thm:cone}.

An earlier accuracy-level audit of this model reported that $99.99\%$ of its $1$D features are exactly zero and that the $0$D and $0$D$+$$1$D banks are equivalent, and inferred from Theorem~\ref{thm:cone}(c) that near-universal coning would account for it --- without checking $\defect_0$ directly. Both halves now hold on the graphs themselves: $1$D persistence is exactly zero on \MistralHOneZero{} of them, and $\defect_0=0$ is verified on \MistralConed{}. The length-normalized rank correlation between $P_0/(N{-}1)$ and $1-\bar m$ has median \MistralRhoNorm{} over layers and never falls below \MistralRhoNormMin{}, so on this model $0$D total persistence is, up to rank, an affine function of sink attention in every layer. Per head, \MistralTohaConed{} of its graphs are exactly prompt-coned, the median $\rho(d,\bar\pi)$ is \MistralTohaRho{}, and Proposition~\ref{prop:toha} has \MistralTohaViol{} violations.
}

\newcommand{\TIMINGPARAGRAPH}{%
Median per-matrix wall-clock time on real attention from TinyLlama-1.1B-Chat-v1.0 (sink features / $\defect_0$ / dense MST / ripser $0$D / ripser $0$D+$1$D) is 0.001/0.02/0.10/0.02/0.02\,ms at $N{=}32$; 0.001/0.03/0.21/0.03/0.03\,ms at $N{=}64$; 0.001/0.07/0.42/0.07/0.06\,ms at $N{=}128$; 0.001/0.25/0.90/0.18/0.18\,ms at $N{=}256$. On i.i.d.\ random matrices of the same size, ripser $0$D+$1$D takes 0.171\,ms ($N{=}32$), 3.46\,ms ($N{=}64$), 115\,ms ($N{=}128$), 4,631\,ms ($N{=}256$). On real attention the same computation is four orders of magnitude cheaper at $N{=}256$, consistent with Theorem~\ref{thm:cone}: sink-dominated graphs have few, short $1$D bars to compute. Sink features, read off one column, are one to two orders of magnitude cheaper still.

}

\newcommand{\RESULTSTOHA}{%
We test Proposition~\ref{prop:toha} on every head of every example: \TohaCells{} head graphs from \TohaSettings{} settings (Table~\ref{tab:toha}). We then run TOHA's head-selection algorithm \citep[Alg.~1]{toha} inside our cross-validation, selecting heads on the whole training fold, which is larger than TOHA's probe set. We run it once on TOHA's score $d$ and, unchanged, on the two first-order scores the proposition names: the $\bar\pi$-score $\frac1{|R|}\sum_{u}(1-\pi_u)$ and the response sink score $1-\bar a_0=\frac1{|R|}\sum_u(1-A_{u0})$.

\emph{The identity holds on real heads.} The sandwich has \TohaViol{} violations. The score $d$ is our own implementation, so we also check it against the authors' released code: on \NativeCells{} head graphs from \NativeSettings{} settings, their \textsc{MTop-Div} routine \citep{toha}, which uses ripser rather than our dense Prim's algorithm, agrees with ours to \NativeMaxDiff{}. Per setting, \TohaConedRange{} of head graphs are exactly prompt-coned, and on those $d$ equals the $\bar\pi$-score up to float rounding (largest gap \TohaGapMax{}). Across examples, the median per-head rank correlation between $d$ and the $\bar\pi$-score is \TohaRhoRange{}. With the sink score it is \TohaRhoSinkRange{} in the sink models, but \TohaRhoSinkNoSink{} in \TohaNoSinkModels{}, which have no first-token sink: there $d$ tracks the largest prompt attention just as tightly while being, if anything, mildly anti-correlated with attention to token~0. The reduction is to max-prompt attention; that this coincides with the sink is a property of the model, not of the theorem. Shifting the sink's attention logit moves these quantities as predicted (Appendix~\ref{app:causal}).

\emph{Topology adds almost nothing to the first-order score.} Stacking $d$ onto the $\bar\pi$-score changes AUC by \TohaLFMaxpRange{}. The change is equivalent to zero within $\pm0.015$ in \TohaLFMaxpEquiv{} of \TohaLFMaxpTotal{} settings, although the 95\% interval lies above $0$ in \TohaLFMaxpPos{} of them. With supervised probes on all heads (the setting of TOHA's own feature comparison), the $\bar\pi$-score differs from $d$ by \TohaSupMaxpRange{} (equivalent in \TohaSupMaxpEquiv{} of \TohaSupMaxpTotal{}). The response sink score differs from $d$ by \TohaSupSinkRange{} (equivalent in \TohaSupSinkEquiv{}). The exceptions are the models without a first-token sink (\TohaNoSinkModels{}), where the sink score is worse, and one on-policy setting whose interval extends past the margin on the positive side. TOHA's discrete head selection is noisier: run on the $\bar\pi$-score it changes AUC by \TohaMaxpRange{} (significantly higher in \TohaMaxpPos{} settings, lower in \TohaMaxpNeg{}, non-inferior at margin $0.015$ in \TohaMaxpNonInf{} of \TohaMaxpTotal{}). Run on the sink score it loses up to \TohaSinkLossSink{} AUC in the sink models (\TohaSinkLossNoSink{} in \TohaNoSinkModels{}, which have no first-token sink): the selected heads are prompt-coned on \TohaSelConedRange{} of examples, but the sink is every response token's top prompt token on only \TohaSelArgZeroRange{}. Stacked onto the non-topological probes, TOHA changes AUC by \TohaLFNTRange{} (95\% interval above $0$ in \TohaLFNTPos{} of \TohaLFNTTotal{}); by the above, whatever it adds is carried by the $\bar\pi$-score.

\emph{Why ``attention to $\langle s\rangle$'' looked weak.} \citet[Table~9]{toha} report that attention to the first token is a much weaker feature than their divergence. Proposition~\ref{prop:toha} identifies the quantity that $d$ actually tracks: per head, the attention of the \emph{response} tokens to their most-attended prompt token. In the sink models, that token is the sink for every response token in \TohaArgZeroRange{} of head graphs. 
}

\newcommand{\CAUSALTEXT}{The intervention leaves the identity intact (no violations in any run). It moves the structure the proposition says matters. Suppressing the sink ($b{=}{-2}$) lowers the share of prompt-coned head graphs from \CausalQwConedZero{} to \CausalQwConedLo{} (Qwen) and from \CausalTlConedZero{} to \CausalTlConedLo{} (TinyLlama). It lowers the share where the sink is every response token's top prompt token from \CausalQwArgZero{} to \CausalQwArgLo{} and from \CausalTlArgZero{} to \CausalTlArgLo{}, and the median rank correlation between $d$ and the response sink score from \CausalQwRhoZero{} to \CausalQwRhoLo{} and from \CausalTlRhoZero{} to \CausalTlRhoLo{}. Strengthening it ($b{=}4$) raises the correlation to \CausalQwRhoHi{} and \CausalTlRhoHi{}. In TinyLlama all three quantities increase monotonically with $b$. In Qwen they peak at $b{=}2$ and fall slightly at $b{=}4$. TOHA's own AUC barely moves (\CausalQwTohaRange{} and \CausalTlTohaRange{}). Its gap to the sink-score version is not monotone in $b$, so the intervention supports the mechanism at the level of the attention graph, not a dose-response in detection accuracy.}
\newcommand{\RESULTSCAUSAL}{%
Proposition~\ref{prop:toha} predicts that TOHA should track the response sink score more closely as the sink strengthens. We test this by intervening on the sink directly. For Qwen2.5-1.5B and TinyLlama on TruthfulQA, we add a constant $b\in\{-2,0,2,4\}$ to the first token's attention logit in every head and recompute all per-head quantities (Table~\ref{tab:causal}). \CAUSALTEXT
}

\newcommand{\RESULTSRECOMMENDATIONS}{%
(i) Report sink attention and $\defect_0$ (or the relevant defect, such as $\defect_P$ for TOHA) next to any attention-graph topological feature, and compare the detector with the first-order statistic that the reduction names.
(ii) Compare topological detectors against per-head non-topological statistics (sink mass, entropy, Lookback ratios) at the same resolution, and against a hidden-state probe.
(iii) Evaluate on the model's own generations; teacher-forced benchmarks mix detection with recognition of dataset authorship.
(iv) If structure beyond the sink is the goal, remove the sink first (deflation) or restrict to answer tokens. Both expose real non-sink geometry, but in our settings it was not new information.
}
